% ============================================================
% Manuscrito completo para revisión con el tutor
% Template: Elsevier CAS Single Column (cas-sc)
% Contenido actual: artículo completo con resultados confirmatorios
% ============================================================

\documentclass[a4paper,fleqn]{cas-sc}

\usepackage[T1]{fontenc}
\usepackage[authoryear,longnamesfirst]{natbib}
\usepackage[spanish,es-tabla]{babel}
\usepackage{amsmath,amsfonts,amssymb}
\usepackage{graphicx}
\usepackage{booktabs}
\usepackage{placeins}
\usepackage{hyperref}

\def\tsc#1{\csdef{#1}{\textsc{\lowercase{#1}}\xspace}}

\begin{document}
\let\WriteBookmarks\relax
\def\floatpagepagefraction{1}
\def\textpagefraction{.001}

% --- Titulo y metadata ---
\shorttitle{Selección y concatenación de descriptores heterogéneos}
\shortauthors{C. Ayala y J. Delgado}

\title[mode=title]{Selección y concatenación de descriptores heterogéneos para la clasificación de texturas}

\tnotemark[1]
\tnotetext[1]{Borrador de trabajo de Tesis de Grado, Facultad Politécnica, Universidad Nacional de Asunción (UNA-FP). Documento preparado para revisión con el tutor; el bloque experimental primario ha finalizado y permanece sujeto a revisión académica.}

% --- Autores ---
\author[1]{Carlos Ayala}
\cormark[1]
\fnmark[1]
\ead{carlos.ayala@est.fpuna.edu.py}
\credit{Conceptualización, Metodología, Software, Experimentación, Análisis formal, Redacción -- borrador original}

\affiliation[1]{organization={Universidad Nacional de Asunción, Facultad Politécnica (UNA-FP)},
            addressline={Campus Universitario},
            city={San Lorenzo},
            state={Central},
            country={Paraguay}}

\author[1]{Joaquín Delgado}
\fnmark[1]
\ead{joaquin.delgado@est.fpuna.edu.py}
\credit{Conceptualización, Metodología, Software, Experimentación, Análisis formal, Redacción -- revisión y edición}

\cortext[1]{Autor de correspondencia}
\fntext[1]{Ambos autores contribuyeron al diseño y ejecución del estudio. Tutor: Prof. José Vázquez.}

% --- Resumen ---
\begin{abstract}
Este trabajo estudia la selección y concatenación de descriptores heterogéneos para clasificación de texturas. Se comparan 20 representaciones clásicas y aprendidas mediante un protocolo anidado que separa la selección del subconjunto de su evaluación externa. El estudio contrasta el mejor descriptor individual, la concatenación completa, Greedy Forward Selection (GFS) y la mejor selección restringida a una familia sobre DTD, FMD y CUReT, con SVM lineal y ResMLP. Las 66 condiciones externas produjeron 264 evaluaciones principales y 6.600 controles con subconjuntos aleatorios. En DTD y FMD, GFS mejora el macro-F1 medio de la concatenación completa; en CUReT presenta diferencias de solo 0,0007 y 0,0004 sobre ocho biparticiones aleatorias, dentro del margen de equivalencia práctica predefinido. Los subconjuntos contienen entre 3,6 y 6,1 extractores en promedio y reducen la dimensionalidad entre 71,5\% y 85,2\%. GFS supera en promedio al mejor individual en las seis combinaciones dataset--clasificador, aunque la composición elegida muestra estabilidad moderada o baja. La evidencia indica que la selección de bloques puede producir representaciones compactas sin pérdida práctica de desempeño, pero no respalda que mezclar familias sea universalmente superior.
\end{abstract}

% --- Highlights ---
\begin{highlights}
\item Se comparan 20 descriptores clásicos y aprendidos mediante evaluación anidada en tres datasets de texturas.
\item GFS mejora o iguala prácticamente la concatenación completa usando entre 71,5\% y 85,2\% menos dimensiones.
\item La utilidad de mezclar familias depende del dataset y no determina una composición universal de descriptores.
\end{highlights}

% --- Palabras clave ---
\begin{keywords}
clasificación de texturas \sep concatenación de descriptores \sep selección de características \sep fusión de representaciones
\end{keywords}

\maketitle

% ============================================================
% 1. INTRODUCCION
% ============================================================
\section{Introducción}\label{sec:introduccion}

La clasificación de texturas es un problema fundamental en visión por computador y constituye un componente relevante en aplicaciones como la inspección de materiales, la recuperación de imágenes por contenido, la teledetección y el análisis automatizado de superficies \citep{bel2022texture}. A diferencia de tareas dominadas por la forma global de los objetos, el reconocimiento de texturas depende de regularidades locales, distribuciones espaciales, respuestas en frecuencia y patrones que pueden variar con la escala, la iluminación y el punto de vista. Esta diversidad explica por qué no existe un descriptor que resulte óptimo para todos los dominios y condiciones de adquisición \citep{cimpoi2014dtd,cimpoi2016fisher}.

Los primeros sistemas de clasificación se apoyaron en descriptores diseñados manualmente. Local Binary Patterns codifica microestructuras locales mediante comparaciones de intensidad \citep{ojala2002lbp}; las matrices de coocurrencia representan relaciones estadísticas entre niveles de gris \citep{haralick1979glcm}; y los histogramas de gradientes describen la distribución local de orientaciones \citep{dalal2005hog}. Estos métodos ofrecen interpretabilidad y bajo costo, pero su capacidad queda condicionada por las propiedades definidas durante su diseño. Las redes convolucionales ampliaron el espacio representacional mediante características aprendidas en grandes conjuntos de imágenes \citep{he2016resnet,tan2019efficientnet}. Posteriormente, los Vision Transformers y los métodos autosupervisados introdujeron representaciones globales y preentrenamientos con menor dependencia de etiquetas \citep{dosovitskiy2021vit,liu2021swin,caron2021dino,oquab2024dinov2}. Como resultado, actualmente conviven descriptores con sesgos inductivos, dimensionalidades y costos muy diferentes.

La coexistencia de estas familias plantea una alternativa a la selección de un único extractor: combinar varias representaciones de la misma imagen. La concatenación de características conserva la información producida por cada descriptor y permite que un clasificador posterior determine su utilidad conjunta. Trabajos previos han mostrado que la fusión de características manuales y profundas puede mejorar la clasificación de texturas \citep{sanchez2013image,cimpoi2016fisher,zhang2017deepten}. Sin embargo, incorporar más vectores no garantiza una representación mejor. Los componentes pueden ser redundantes, operar en escalas incompatibles o incrementar la dimensionalidad hasta perjudicar la generalización y el costo computacional. Por ello, el problema no consiste solamente en decidir si se deben concatenar descriptores, sino en determinar cuáles aportan información complementaria y cuándo conviene incorporarlos.

Una parte de la literatura evalúa combinaciones pequeñas o concatenaciones definidas de antemano. Este enfoque permite demostrar que dos representaciones pueden complementarse, pero ofrece información limitada sobre el comportamiento de espacios más amplios y heterogéneos. También dificulta separar tres efectos: la calidad de los descriptores individuales, el beneficio de aumentar la dimensionalidad y la utilidad específica de seleccionar un subconjunto. Esta distinción es especialmente importante cuando se comparan descriptores clásicos, redes convolucionales, Transformers y modelos autosupervisados bajo diferentes clasificadores y tamaños de muestra.

El presente trabajo estudia la siguiente pregunta: \emph{¿en qué condiciones la concatenación de descriptores heterogéneos mejora la clasificación de texturas, y qué balance ofrece la selección de subconjuntos entre desempeño y dimensionalidad?} Para responderla se evaluaron 20 extractores distribuidos en cuatro familias: cinco descriptores clásicos, seis redes convolucionales, cuatro Vision Transformers y cinco representaciones con preentrenamiento autosupervisado o multimodal. La etapa exploratoria abarcó seis colecciones. Después de auditar la procedencia, el orden de las muestras y la independencia de las particiones, el bloque confirmatorio quedó restringido a DTD, FMD y CUReT, con SVM lineal y ResMLP como clasificadores primarios. Esta reducción obedeció a criterios de integridad de datos y no al rendimiento observado.

El diseño compara cuatro configuraciones. La primera selecciona el mejor descriptor individual como referencia. La segunda concatena los 20 extractores en una representación fija. La tercera emplea Greedy Forward Selection (GFS) para construir un subconjunto incremental a partir del macro-F1 interno. La cuarta restringe la misma búsqueda a una familia representacional. Esta comparación permite separar el efecto de la dimensionalidad, la selección de bloques y la heterogeneidad, además de analizar la dependencia de la solución respecto del dataset y del clasificador.

Los resultados confirmatorios muestran un patrón condicionado por el dominio. GFS mejora el macro-F1 medio de la concatenación completa en DTD y FMD; en CUReT queda entre 0,0004 y 0,0007 por debajo sobre ocho biparticiones aleatorias, una diferencia dentro del margen de equivalencia práctica de 0,01, mientras utiliza entre un quinto y un cuarto de las dimensiones. Frente al mejor descriptor individual, GFS obtiene una diferencia media positiva en las seis combinaciones dataset--clasificador. La composición seleccionada cambia entre particiones y presenta estabilidad moderada o baja, por lo que la complementariedad no puede reducirse a una combinación universal.

La contribución del trabajo es principalmente empírica y metodológica. En primer lugar, se ofrece una comparación homogénea de representaciones visuales heterogéneas sobre varios datasets y clasificadores. En segundo lugar, se distingue el efecto de concatenar todas las características del efecto de seleccionar componentes complementarios. En tercer lugar, se cuantifica el compromiso entre macro-F1, número de extractores y dimensionalidad. Finalmente, se documentan configuraciones en las que la concatenación no mejora, lo cual permite delimitar las condiciones de utilidad del enfoque y evita interpretar la fusión como una operación universalmente beneficiosa.

El resto del artículo se organiza de la siguiente manera. La segunda sección revisa los descriptores de texturas y las estrategias de fusión relacionadas. La tercera describe los datasets, extractores, clasificadores y protocolos de selección y evaluación. La cuarta presenta los resultados confirmatorios. La quinta analiza la complementariedad, el costo dimensional y las amenazas a la validez. Finalmente, la sexta sección resume las conclusiones del estudio.

% ============================================================
% 2. TRABAJOS RELACIONADOS
% ============================================================
\section{Trabajos relacionados}\label{sec:trabajos}

\subsection{Representaciones diseñadas y aprendidas para texturas}

La descripción computacional de texturas se ha abordado históricamente mediante representaciones que resumen regularidades locales sin requerir aprendizaje supervisado. Las matrices de coocurrencia caracterizan dependencias estadísticas entre intensidades \citep{haralick1979glcm}; LBP codifica configuraciones binarias en vecindarios locales y admite extensiones invariantes a rotación y escala \citep{ojala2002lbp}; HOG agrega orientaciones de gradiente en regiones espaciales \citep{dalal2005hog}; y DRLBP refuerza la invariancia frente a cambios de orientación \citep{mehta2011drlbp}. Estos descriptores expresan propiedades diferentes de una misma imagen, por lo que no son intercambiables: una representación basada en microestructuras no conserva necesariamente la misma información que otra basada en coocurrencias, frecuencia u orientación.

Las representaciones aprendidas trasladaron parte de esa decisión de diseño al proceso de preentrenamiento. Las redes convolucionales jerárquicas, como ResNet y EfficientNet, producen activaciones que pueden reutilizarse como vectores de características \citep{he2016resnet,tan2019efficientnet}. En reconocimiento de texturas, la agregación de respuestas convolucionales mostró que modelos entrenados originalmente para objetos también capturan estadísticas locales útiles \citep{cimpoi2016fisher}. Los Vision Transformers ampliaron este repertorio mediante interacciones de largo alcance entre regiones de la imagen \citep{dosovitskiy2021vit,liu2021swin}. A su vez, los objetivos autosupervisados y multimodales generan espacios de representación sin depender exclusivamente de las etiquetas del problema de destino \citep{caron2021dino,oquab2024dinov2,he2022mae,zhai2023siglip}. En este trabajo, estas arquitecturas no se tratan como fines en sí mismas, sino como fuentes heterogéneas de descriptores que pueden contener información coincidente o complementaria.

\subsection{Fusión y concatenación de características}

La fusión de características combina información antes de la decisión final del clasificador. Una estrategia directa consiste en concatenar los vectores extraídos de una misma imagen. A diferencia de la fusión de decisiones, esta operación conserva explícitamente las coordenadas de cada representación y permite entrenar un único clasificador sobre el espacio conjunto. Su sencillez constituye una ventaja experimental: la mejora puede atribuirse a la interacción entre descriptores sin introducir un módulo de fusión altamente parametrizado. No obstante, el vector resultante crece con cada extractor incorporado y puede acumular escalas incompatibles, variables redundantes y dimensiones poco informativas.

Los Fisher vectors mostraron que la agregación de descriptores locales puede producir representaciones de alto rendimiento para clasificación de imágenes \citep{sanchez2013image}. En texturas, los bancos de filtros profundos combinaron activaciones convolucionales con esquemas de agregación estadística y establecieron una relación efectiva entre características aprendidas y modelos clásicos de orden local \citep{cimpoi2016fisher}. DeepTEN integró una capa de codificación residual para recuperar priors estadísticos propios de las texturas dentro de una arquitectura profunda \citep{zhang2017deepten}. Estos trabajos respaldan la utilidad de combinar señales, pero estudian mecanismos de agregación o arquitecturas específicas. El problema considerado aquí es distinto: dados múltiples extractores ya disponibles y congelados, se busca determinar qué subconjunto conviene concatenar bajo un presupuesto dimensional y un clasificador común.

La concatenación completa proporciona una referencia importante, aunque no constituye por sí sola un procedimiento de selección. Si todos los descriptores se incorporan de manera automática, una mejora puede deberse a una única representación dominante, mientras que el resto solo aumenta el costo. De forma análoga, una concatenación acumulativa en orden fijo mezcla la utilidad del descriptor añadido con su posición en la secuencia. Para separar estos efectos es necesario comparar, bajo las mismas particiones, al mejor descriptor individual, la concatenación de todos los componentes y una búsqueda explícita de subconjuntos.

\subsection{Selección de subconjuntos y vacío experimental}

La selección de características suele operar sobre coordenadas individuales. En cambio, cada candidato de este estudio es un bloque completo generado por un extractor, de modo que la unidad de selección conserva su identidad, dimensión y familia representacional. Greedy Forward Selection (GFS) resulta adecuado para este escenario porque construye el conjunto de forma incremental: en cada paso evalúa qué bloque no seleccionado produce la mayor mejora al añadirse a la combinación vigente. El procedimiento no garantiza el óptimo global, pero hace tratable un espacio de búsqueda combinatorio y genera una trayectoria interpretable entre desempeño y costo.

La evidencia disponible deja tres aspectos que motivan esta investigación. Primero, las comparaciones suelen concentrarse en una representación o en combinaciones diseñadas previamente, por lo que describen de manera incompleta la redundancia entre familias heterogéneas. Segundo, no siempre se contrasta la selección con subconjuntos aleatorios del mismo tamaño, requisito necesario para distinguir una elección informada de una ganancia atribuible únicamente a concatenar más dimensiones. Tercero, seleccionar y evaluar sobre las mismas particiones puede producir estimaciones optimistas. Por ello, el protocolo de este trabajo separa la etapa exploratoria de una evaluación confirmatoria anidada: la selección se realiza exclusivamente dentro de los datos de entrenamiento de cada partición externa y el desempeño se estima sobre observaciones no utilizadas para escoger el subconjunto.

En síntesis, la literatura demuestra que tanto los descriptores diseñados como los aprendidos pueden representar información relevante de textura, y que la agregación de características puede mejorar la discriminación. La contribución buscada no es proponer otro extractor ni atribuir el resultado a una arquitectura particular. El foco está en evaluar sistemáticamente la \emph{concatenación de descriptores como problema de selección de bloques}, cuantificando cuándo la complementariedad justifica el aumento de dimensionalidad y cuándo una representación más compacta ofrece una solución equivalente o superior.

% ============================================================
% 3. METODOLOGIA
% ============================================================
\section{Metodología}\label{sec:metodologia}

\subsection{Diseño del estudio}

El estudio evalúa la concatenación como un problema de selección de bloques. Cada extractor transforma una imagen $x_i$ en un vector $\mathbf{z}_{ij}\in\mathbb{R}^{d_j}$, donde $j$ identifica el descriptor y $d_j$ su dimensionalidad. Para un subconjunto $S$ de extractores, la representación conjunta se define como
\begin{equation}
 \mathbf{z}_{i,S}=\mathop{\Vert}_{j\in S}\frac{\mathbf{z}_{ij}}{\max(\lVert\mathbf{z}_{ij}\rVert_2,10^{-12})},
 \label{eq:concat}
\end{equation}
donde $\Vert$ denota concatenación. La normalización $L_2$ se aplica por descriptor y por muestra antes de unir los bloques, para evitar que la escala de un extractor determine por sí sola la contribución al clasificador.

La pregunta primaria es si GFS mantiene o mejora el macro-F1 del mejor descriptor individual y de la concatenación completa cuando la selección se realiza sin consultar la partición externa de evaluación. Se considera además el costo de la solución mediante el número de extractores $k$, la dimensionalidad total y el tiempo de ajuste. El margen de equivalencia práctica se fijó antes del análisis en $|\Delta\text{macro-F1}|<0{,}01$.

\subsection{Datasets y control de integridad}

La etapa exploratoria incluyó DTD, FMD, Outex13, CUReT, Soil y VisTex. Antes del análisis confirmatorio se construyó un manifiesto canónico por dataset y se verificaron etiquetas, orden, dimensiones, valores finitos, hashes de las imágenes y grupos de dependencia. La Tabla~\ref{tab:datasets} resume la decisión resultante. DTD, FMD y CUReT conforman el bloque principal; Outex se incorporó posteriormente como extensión sobre su colección y partición oficiales.

\FloatBarrier
\begin{table}[!ht]
\centering
\caption{Datasets considerados y decisión del control de integridad.}
\label{tab:datasets}
\small
\resizebox{\textwidth}{!}{%
\begin{tabular}{lrrl}
\toprule
Dataset & Muestras & Clases & Protocolo confirmatorio \\
\midrule
DTD & 5.640 & 47 & 10 splits oficiales \\
FMD & 1.000 & 10 & 3 semillas $\times$ 5 folds agrupados \\
CUReT & 5.612 & 61 & 8 biparticiones aleatorias de 46 condiciones \\
Outex13 oficial & 1.360 & 68 & 680 entrenamiento / 680 test \\
Soil & 1.186 & -- & excluido: archivos no legibles \\
VisTex & 167 & 19 oficiales & excluido: unidades independientes insuficientes \\
\bottomrule
\end{tabular}}
\end{table}
\FloatBarrier

DTD contiene 5.640 imágenes de texturas en condiciones no controladas y proporciona diez particiones oficiales de entrenamiento, validación y prueba \citep{cimpoi2014dtd}. En cada partición se unieron entrenamiento y validación para ejecutar la selección interna, mientras que el test oficial se reservó para una única evaluación. Los duplicados byte-idénticos que cruzaban roles fueron retirados del conjunto de entrenamiento, conservando intacta la evaluación.

FMD reúne 1.000 fotografías de diez categorías de materiales \citep{sharan2013fmd}. Los 20 descriptores se reextrajeron desde las imágenes oficiales en un orden canónico. Cada imagen presentó un hash único y se utilizó como unidad de agrupación. La evaluación externa empleó cinco folds estratificados y agrupados, repetidos con semillas 42, 123 y 2026.

Para CUReT se utilizó la distribución gris recortada de 61 materiales, con 92 condiciones compartidas por clase \citep{varma2003filter}. Cada condición constituye un grupo para impedir que la misma configuración de adquisición aparezca a ambos lados de una partición. La evaluación externa emplea ocho biparticiones aleatorias e independientes de las 92 condiciones en mitades de 46, generadas con semillas fijas 1 a 8. Como cada condición contiene exactamente una imagen por clase, toda división de grupos completos queda balanceada entre las 61 clases sin necesidad de estratificar por separado. Este esquema conserva el tamaño 46/46 publicado y, al no reutilizar una única asignación, evita que las conclusiones dependan de un solo corte. No se presenta como el split histórico exacto.

Una reproducción determinista alternada de dos mitades complementarias, empleada en una etapa previa del estudio, arrojó resultados concordantes con los aquí reportados. No se promedia con las biparticiones aleatorias porque sus dos direcciones son la misma partición invertida y, por tanto, no constituyen réplicas independientes.

Soil y VisTex no se descartaron por bajo desempeño: Soil contenía archivos vacíos o ilegibles y VisTex no disponía de suficientes unidades fuente independientes para combinar validación externa e interna. Outex se restauró desde la distribución completa y se evalúa únicamente con su partición oficial; no se utilizan resultados exploratorios anteriores.

\subsection{Espacio de descriptores}

El espacio de candidatos contiene 20 extractores congelados agrupados por su mecanismo representacional:
\begin{itemize}
 \item \textbf{clásicos}: LBP, DRLBP, Gabor, GLCM y HOG;
 \item \textbf{CNN}: VGG16, ResNet-50, ResNet-101, DenseNet-121, EfficientNet-B0 y ConvNeXt V2-T;
 \item \textbf{Transformers}: ViT-B/16, DeiT-S, Swin-T y EVA-02 base;
 \item \textbf{preentrenamiento autosupervisado o multimodal}: DINOv2 small, base y large, MAE base y SigLIP base.
\end{itemize}
Los vectores concatenados de los 20 extractores suman 19.628 dimensiones. Todos fueron calculados una sola vez y almacenados junto con sus etiquetas y metadatos. El estudio no ajusta los backbones ni compara estrategias de \emph{fine-tuning}; evalúa exclusivamente la utilidad de seleccionar y concatenar representaciones fijas.

\subsection{Clasificadores y comparadores}

Se emplearon dos clasificadores primarios con sesgos diferentes. La SVM lineal utiliza $C=1$, ponderación balanceada por clase y un máximo de 10.000 iteraciones. ResMLP proyecta el vector de entrada a 256 unidades y aplica tres bloques residuales, \emph{dropout} de 0,1 y una capa de salida multiclase. Se entrena con AdamW, tasa de aprendizaje $10^{-3}$, decaimiento de pesos $10^{-4}$, lotes de 256 muestras, un máximo de 100 épocas y parada temprana con paciencia 10. Las semillas se fijan para NumPy y PyTorch; cuando está disponible, ResMLP se ejecuta en GPU.

Para cada partición externa se comparan cuatro estrategias:
\begin{enumerate}
 \item \textbf{mejor individual}: el extractor con mayor macro-F1 interno;
 \item \textbf{concatenación completa}: los 20 extractores en una representación fija;
 \item \textbf{GFS}: el mejor subconjunto construido por selección voraz;
 \item \textbf{mejor familia homogénea}: GFS restringido a una familia y con un máximo de componentes no mayor que el subconjunto heterogéneo.
\end{enumerate}
El cuarto comparador permite distinguir si una ganancia procede de combinar familias o simplemente de seleccionar varios extractores similares.

\subsection{Greedy Forward Selection}

GFS comienza con un conjunto vacío. En el paso $t$, evalúa cada descriptor aún no seleccionado al añadirlo al subconjunto vigente y elige el candidato con mayor macro-F1 medio en la validación interna:
\begin{equation}
 j_t=\arg\max_{j\notin S_{t-1}}\widehat{F}_{1,\mathrm{inner}}(S_{t-1}\cup\{j\}),
 \qquad S_t=S_{t-1}\cup\{j_t\}.
 \label{eq:gfs}
\end{equation}
La búsqueda se limita a $k\leq 8$. Se detiene cuando la mejora respecto del paso anterior es menor que 0,001 durante dos pasos consecutivos, pero devuelve el mejor paso histórico, no necesariamente el último. Los empates se resuelven de forma determinista mediante el nombre del extractor. Todas las evaluaciones candidatas utilizan únicamente muestras pertenecientes al entrenamiento externo.

\subsection{Protocolo anidado y controles aleatorios}

La selección interna utiliza cuatro folds estratificados y agrupados. En DTD se ejecuta dentro de la unión \emph{train+val} de cada split oficial; en FMD dentro del entrenamiento de cada uno de los 15 folds externos; y en CUReT dentro de la mitad de 46 condiciones asignada al entrenamiento. Después de seleccionar el subconjunto y los comparadores, cada modelo se ajusta con todo el entrenamiento externo y se evalúa una sola vez en el test correspondiente. Se verifica programáticamente que ningún grupo aparezca en entrenamiento y prueba.

Como control, para cada condición externa se generan 100 subconjuntos aleatorios sin reemplazo con el mismo $k$ que GFS. El percentil del subconjunto elegido y su valor $p$ empírico se calculan como
\begin{equation}
 p_{\mathrm{emp}}=\frac{1+\#\{F_{1,\mathrm{random}}\geq F_{1,\mathrm{GFS}}\}}{B+1},\qquad B=100.
 \label{eq:pemp}
\end{equation}
Este control responde si la selección informada mejora lo esperable al concatenar, al azar, la misma cantidad de bloques.

\subsection{Métricas y análisis}

La métrica primaria es macro-F1 en el test externo; la exactitud se conserva como diagnóstico secundario. También se registran $k$, dimensionalidad, tiempo de ajuste, frecuencia de selección por extractor y familia, y estabilidad de Jaccard entre subconjuntos. Los contrastes primarios son GFS frente al mejor individual y GFS frente a la concatenación completa.

Para los tres datasets se reportan medias, desviaciones estándar, diferencias pareadas, conteos de victorias/empates/derrotas e intervalos bootstrap pareados. Estos intervalos se consideran descriptivos, porque los splits oficiales de DTD, las repeticiones de FMD y las biparticiones de CUReT reutilizan las mismas imágenes y no constituyen réplicas completamente independientes; no se calculan pruebas de hipótesis formales sobre estas diferencias por la misma razón. Una diferencia pequeña se interpreta mediante el margen práctico predefinido y no mediante la ausencia de significación estadística.

% ============================================================
% 4. RESULTADOS
% ============================================================
\section{Resultados}\label{sec:resultados}

\subsection{Alcance de la evidencia disponible}

Finalizaron y superaron los controles de integridad las 20 condiciones de DTD (10 splits por dos clasificadores), las 30 de FMD (15 folds por dos clasificadores) y las 16 de CUReT (ocho biparticiones por dos clasificadores). En conjunto produjeron 264 evaluaciones principales y 6.600 controles aleatorios. No se detectaron filas duplicadas, condiciones faltantes ni valores no finitos en las métricas obligatorias. Como comprobación adicional, una repetición determinista reducida de DTD--SVM y una repetición completa de FMD--ResMLP reprodujeron exactamente las métricas, los subconjuntos seleccionados, $k$ y la dimensionalidad de sus ejecuciones de referencia.

\subsection{Comparación de estrategias}

La Tabla~\ref{tab:resultados-confirmatorios} presenta el macro-F1 externo medio. GFS obtuvo el mayor promedio frente al mejor individual en las seis combinaciones dataset--clasificador. Frente a concatenar los 20 extractores, la diferencia fue positiva en DTD (0,0046 con SVM y 0,0066 con ResMLP) y FMD (0,0122 y 0,0208), y ligeramente negativa en CUReT ($-0,0007$ y $-0,0004$). Estas dos últimas diferencias se encuentran dentro del margen de equivalencia práctica predefinido.

\FloatBarrier
\begin{table}[!ht]
\centering
\caption{Resultados confirmatorios. Se reporta macro-F1 medio $\pm$ desviación estándar sobre el test externo. $k$ y dimensiones corresponden al promedio de GFS.}
\label{tab:resultados-confirmatorios}
\small
\resizebox{\textwidth}{!}{%
\begin{tabular}{llccccc}
\toprule
Dataset & Clasificador & Individual & Completa & GFS & $k_{\mathrm{GFS}}$ (promedio) & Dim. GFS (promedio) \\
\midrule
DTD & SVM & $0{,}8426\pm0{,}0109$ & $0{,}8648\pm0{,}0076$ & $\mathbf{0{,}8694}\pm0{,}0096$ & 6,1 & 5.590 \\
DTD & ResMLP & $0{,}8373\pm0{,}0089$ & $0{,}8602\pm0{,}0086$ & $\mathbf{0{,}8668}\pm0{,}0094$ & 5,9 & 4.792 \\
FMD & SVM & $0{,}9726\pm0{,}0150$ & $0{,}9644\pm0{,}0167$ & $\mathbf{0{,}9766}\pm0{,}0135$ & 3,6 & 2.909 \\
FMD & ResMLP & $0{,}9690\pm0{,}0133$ & $0{,}9575\pm0{,}0180$ & $\mathbf{0{,}9783}\pm0{,}0123$ & 4,4 & 3.077 \\
CUReT & SVM & $0{,}9942\pm0{,}0021$ & $\mathbf{0{,}9989}\pm0{,}0008$ & $0{,}9982\pm0{,}0010$ & 4,6 & 5.497 \\
CUReT & ResMLP & $0{,}9967\pm0{,}0006$ & $\mathbf{0{,}9988}\pm0{,}0008$ & $0{,}9983\pm0{,}0008$ & 3,8 & 3.408 \\
\bottomrule
\end{tabular}}
\end{table}

La identidad del descriptor reportado como mejor individual no es fija entre particiones. En DTD, SigLIP-base fue el mejor individual en los 10 splits con ResMLP y en 8 de 10 con SVM; los 2 splits restantes correspondieron a DINOv2-large. En FMD, SigLIP-base resultó el mejor individual en los 15 folds con ambos clasificadores. En CUReT, DINOv2-large fue el mejor individual en las ocho biparticiones con SVM, mientras que con ResMLP alternó entre DINOv2-large y DINOv2-base, cuatro veces cada uno. Este patrón muestra que los modelos con preentrenamiento autosupervisado o multimodal más recientes (SigLIP, DINOv2) dominan de forma sistemática como referencia individual, lo cual contextualiza la magnitud de la mejora de GFS: la comparación se da contra un descriptor ya fuerte, no contra un extractor débil.

En DTD, GFS superó al mejor individual en los diez splits con ambos clasificadores. Frente a la concatenación completa obtuvo 8 victorias de 10 con SVM y 9 de 10 con ResMLP. En FMD, GFS superó al mejor individual en 9 de 15 folds con SVM (dos empates) y en 13 de 15 con ResMLP (un empate); frente a la concatenación completa registró 12 y 14 victorias, respectivamente. Estos conteos muestran que la mejora media no depende de una única partición excepcional.

Los intervalos bootstrap pareados del cambio de GFS frente al mejor individual fueron $[0{,}0226,0{,}0314]$ y $[0{,}0244,0{,}0339]$ en DTD con SVM y ResMLP, respectivamente. En FMD fueron $[-0{,}0014,0{,}0090]$ y $[0{,}0062,0{,}0123]$: el intervalo de SVM cruza cero y aconseja interpretar esa diferencia como incierta y pequeña. Frente a la concatenación completa, los intervalos fueron positivos en DTD ($[0{,}0007,0{,}0081]$ y $[0{,}0035,0{,}0098]$) y FMD ($[0{,}0068,0{,}0170]$ y $[0{,}0143,0{,}0271]$). Debido al solapamiento entre splits, estos intervalos describen la variación observada y no se interpretan como inferencia sobre réplicas independientes.

En CUReT, GFS superó al mejor individual en las ocho biparticiones con ambos clasificadores, con intervalos bootstrap enteramente positivos ($[0{,}0028,0{,}0052]$ con SVM y $[0{,}0012,0{,}0021]$ con ResMLP). Frente a la concatenación completa el balance se invierte y se estrecha: 2 victorias, 1 empate y 5 derrotas con SVM, y 3 victorias, 1 empate y 4 derrotas con ResMLP, con intervalos que cruzan cero ($[-0{,}0014,0{,}0000]$ y $[-0{,}0011,0{,}0002]$). La totalidad de las 16 diferencias externas quedó por debajo del margen de 0,01, por lo que el resultado se interpreta como equivalencia práctica y no como superioridad de la representación completa.

La concatenación completa contiene 19.628 dimensiones. Los subconjuntos GFS redujeron ese tamaño en 71,5\%--75,6\% en DTD, 84,3\%--85,2\% en FMD y 72,0\%--82,6\% en CUReT. El resultado global no es que una representación más grande sea siempre mejor, sino que la selección interna puede conservar una fracción del espacio y mejorar o mantener el desempeño dentro de un margen práctico.

La reducción dimensional también disminuyó el tiempo medio de ajuste. Con SVM, GFS requirió 13,12 frente a 58,74 segundos en DTD, 0,78 frente a 4,05 en FMD y 11,67 frente a 57,51 en CUReT. Con ResMLP, los tiempos fueron 3,96 frente a 5,17 segundos, 1,51 frente a 2,24 y 3,07 frente a 4,77, respectivamente. Estas mediciones cubren únicamente el ajuste del clasificador y no el costo de ejecutar los extractores.

\subsection{Heterogeneidad frente a selección homogénea}

La mejor familia homogénea alcanzó macro-F1 de 0,8606 y 0,8611 en DTD con SVM y ResMLP, por debajo de GFS en ambos casos. En FMD, las diferencias fueron menores: con ResMLP, GFS obtuvo 0,9783 frente a 0,9773; con SVM, la selección homogénea alcanzó 0,9789 y superó ligeramente a GFS (0,9766). En CUReT el orden depende del clasificador: la selección homogénea alcanzó 0,9983 frente a 0,9982 de GFS con SVM, mientras que con ResMLP GFS quedó por encima, 0,9983 frente a 0,9979. Estos casos impiden afirmar que mezclar familias sea universalmente superior. La hipótesis compatible con la evidencia es condicional: la heterogeneidad resulta útil cuando aporta información complementaria, pero una familia puede bastar en dominios con desempeño saturado.

\subsection{Composición de los subconjuntos seleccionados}

Para caracterizar qué extractores sostienen el desempeño de GFS, la Tabla~\ref{tab:seleccion-frecuencia} reporta los descriptores que aparecieron en al menos la mitad de las particiones externas. En DTD y FMD, DINOv2-large y SigLIP-base forman un núcleo estable presente en el 100\,\% de las particiones con los dos clasificadores; EVA-02 base acompaña con frecuencia alta en DTD (90\,\% con ResMLP, 100\,\% con SVM) y moderada en FMD (60\,\% con ResMLP). El resto de los extractores --clásicos y CNN-- aparece de forma ocasional (50\,\% o menos), lo que indica que su contribución depende de la partición específica y no de una complementariedad sistemática.

\FloatBarrier
\begin{table}[!ht]
\centering
\caption{Extractores seleccionados por GFS con frecuencia $\geq$ 50\,\% de las particiones externas.}
\label{tab:seleccion-frecuencia}
\small
\begin{tabular}{llcc}
\toprule
Dataset & Clasificador & Extractor & Frecuencia \\
\midrule
DTD & SVM & dinov2\_large & 1,00 (10/10) \\
DTD & SVM & eva02\_base & 1,00 (10/10) \\
DTD & SVM & siglip\_base & 1,00 (10/10) \\
DTD & SVM & dinov2 & 0,50 (5/10) \\
DTD & ResMLP & dinov2\_large & 1,00 (10/10) \\
DTD & ResMLP & siglip\_base & 1,00 (10/10) \\
DTD & ResMLP & eva02\_base & 0,90 (9/10) \\
DTD & ResMLP & densenet121 & 0,50 (5/10) \\
DTD & ResMLP & resnet50 & 0,50 (5/10) \\
FMD & SVM & siglip\_base & 1,00 (15/15) \\
FMD & SVM & dinov2\_large & 0,87 (13/15) \\
FMD & ResMLP & dinov2\_large & 1,00 (15/15) \\
FMD & ResMLP & siglip\_base & 1,00 (15/15) \\
FMD & ResMLP & eva02\_base & 0,60 (9/15) \\
CUReT & SVM & dinov2\_large & 1,00 (8/8) \\
CUReT & SVM & dinov2\_small & 0,75 (6/8) \\
CUReT & SVM & dinov2 & 0,63 (5/8) \\
CUReT & ResMLP & dinov2\_large & 0,75 (6/8) \\
CUReT & ResMLP & dinov2 & 0,50 (4/8) \\
CUReT & ResMLP & dinov2\_small & 0,50 (4/8) \\
\bottomrule
\end{tabular}
\end{table}
\FloatBarrier

CUReT presenta un patrón distinto: los seis descriptores que superan el umbral pertenecen todos a la familia DINOv2, y ningún extractor clásico, convolucional o de otra familia autosupervisada alcanza el 50\,\%. Este comportamiento concuerda con que la selección homogénea resulte competitiva en este dataset y sugiere que, en condiciones de adquisición controladas y con desempeño próximo al techo, la complementariedad entre familias aporta poco.

\subsection{Controles aleatorios y estabilidad}

GFS se situó en percentiles altos de los subconjuntos aleatorios de igual tamaño. El patrón fue más claro en DTD, con percentiles medianos de 0,995 para SVM y 0,980 para ResMLP; en FMD la mediana fue 0,970 con ambos clasificadores. En CUReT, con valores cercanos al techo, los percentiles medianos fueron 0,860 con SVM y 0,815 con ResMLP, y $p_{\mathrm{emp}}$ resultó menor que 0,05 en solo 3 de las 16 biparticiones. Esta menor separación es coherente con un dominio saturado, donde casi cualquier subconjunto suficientemente grande alcanza un desempeño próximo al máximo. Estos valores se reportan como controles por condición y no como una prueba global de significación.

La similitud de Jaccard media entre subconjuntos fue 0,435--0,446 en DTD y 0,400--0,447 en FMD. En CUReT fue menor, 0,337 con SVM y 0,229 con ResMLP sobre 28 pares de biparticiones. Por tanto, un desempeño externo similar no implica que GFS recupere una única composición. La estabilidad observada corresponde mejor a la utilidad predictiva y al tamaño del subconjunto que a la identidad exacta de todos sus componentes.

% ============================================================
% 5. DISCUSION
% ============================================================
\FloatBarrier
\section{Discusión}\label{sec:discusion}

Los resultados apoyan la formulación central de la tesis: la concatenación debe analizarse junto con la selección de componentes. Usar los 20 extractores produjo el mayor promedio en CUReT, pero no en DTD ni FMD y, aun en CUReT, su ventaja sobre GFS fue menor que 0,001 y ninguna de las 16 diferencias externas superó el margen práctico. GFS empleó solo 15\%--29\% de la dimensionalidad completa y mejoró o mantuvo el desempeño dentro del margen práctico. En consecuencia, el aporte no reside en sustituir un descriptor por otro, sino en establecer un procedimiento reproducible para decidir qué bloques conservar.

La magnitud del beneficio depende del dataset. En DTD, el mejor descriptor individual quedó claramente por debajo de las representaciones combinadas, lo que sugiere que las texturas no controladas contienen señales distribuidas entre varios extractores. En FMD, el mejor individual ya alcanzó valores altos y el margen adicional fue menor, especialmente con SVM. Aun así, la concatenación completa rindió peor que GFS, señal de que añadir características redundantes puede perjudicar incluso cuando los componentes individuales son fuertes.

El comparador homogéneo añade un matiz relevante. GFS heterogéneo fue superior en DTD y prácticamente equivalente al mejor conjunto homogéneo en FMD y CUReT; en este último, además, los subconjuntos seleccionados se concentraron casi exclusivamente en la familia autosupervisada. Por ello, la diversidad arquitectónica no produce una ventaja automática. Una interpretación más rigurosa es que ampliar el conjunto de candidatos permite encontrar complementariedad cuando existe; el protocolo de selección, no la presencia de una arquitectura específica, determina si esa diversidad se conserva.

La concordancia cualitativa entre SVM y ResMLP también fortalece el hallazgo. Ambos favorecieron GFS sobre la concatenación completa en DTD y FMD, y ambos mostraron equivalencia práctica entre esas estrategias en CUReT. Las diferencias y los subconjuntos elegidos no fueron idénticos, de modo que el patrón no parece exclusivamente una propiedad de un cabezal lineal o de una red residual. Sin embargo, esto no implica equivalencia general entre clasificadores: cada uno conserva una dinámica de optimización y un costo distintos.

Desde una perspectiva práctica, el ahorro dimensional disminuye almacenamiento, transferencia de características y costo de ajuste. Los tiempos registrados confirman que el ajuste sobre GFS es menor que sobre la concatenación completa en las configuraciones de mayor dimensión. No obstante, $k$, dimensión y tiempo de ajuste no reflejan todo el costo: todavía es necesario ejecutar cada extractor incluido, y dos vectores de igual dimensión pueden requerir tiempos de inferencia diferentes. La ausencia de una medición homogénea del costo de extracción limita cualquier conclusión sobre latencia de extremo a extremo.

\subsection{Comparación con cifras publicadas}

Para situar estos resultados frente a la literatura, la Tabla~\ref{tab:sota} contrasta el macro-F1 de GFS con cifras de accuracy publicadas en trabajos que introdujeron o popularizaron estos bancos de datos. En DTD, \citet{cimpoi2016fisher} reportan 74,7\,\%~$\pm$~1,0 de accuracy con FV-CNN sobre VGG-VD, frente al 86,94\,\% de macro-F1 que alcanza GFS con SVM. En FMD, el mismo trabajo reporta 82,4\,\%~$\pm$~1,5 de accuracy, frente al 97,66\,\%--97,83\,\% de macro-F1 de GFS. En CUReT, \citet{varma2003filter} reportan una accuracy media de 97,28\,\%~$\pm$~0,32 con el banco de filtros MR8 y 46 imágenes de entrenamiento por clase --protocolo cercano al utilizado en este trabajo--, frente al 99,80\,\%--99,88\,\% de macro-F1 de GFS.

Estas comparaciones deben leerse con cautela: las cifras de referencia reportan accuracy y no macro-F1, provienen de protocolos de partición distintos y anteceden en 10 a 20 años a los descriptores autosupervisados y multimodales evaluados en este trabajo. La comparación no aísla el aporte metodológico de GFS del salto que ya proveen los descriptores modernos por sí solos. Aun así, la magnitud de la diferencia es consistente con la mejora esperable al incorporar representaciones más recientes, y sitúa el desempeño absoluto de GFS por encima de los puntos de referencia disponibles en la literatura citada.

\FloatBarrier
\begin{table}[!ht]
\centering
\caption{Comparación con cifras publicadas. No son directamente equiparables: distinta métrica (accuracy vs.\ macro-F1) y protocolo.}
\label{tab:sota}
\small
\resizebox{\textwidth}{!}{%
\begin{tabular}{lllcc}
\toprule
Dataset & Referencia & Método & Cifra publicada & GFS (este trabajo) \\
\midrule
DTD & \citet{cimpoi2016fisher} & FV-CNN (VGG-VD) & 74,7\,\% $\pm$ 1,0 (accuracy) & 86,94\,\% macro-F1 (SVM) \\
FMD & \citet{cimpoi2016fisher} & FV-CNN (VGG-VD) & 82,4\,\% $\pm$ 1,5 (accuracy) & 97,66\,\%--97,83\,\% macro-F1 \\
CUReT & \citet{varma2003filter} & MR8 (46 train/clase) & 97,28\,\% $\pm$ 0,32 (accuracy) & 99,82\,\%--99,83\,\% macro-F1 \\
\bottomrule
\end{tabular}}
\end{table}
\FloatBarrier

\subsection{Limitaciones y amenazas a la validez}

El estudio presenta cinco limitaciones principales. Primero, GFS es una heurística local y no garantiza el subconjunto globalmente óptimo. Segundo, se utilizan extractores congelados y no se evalúa si el ajuste conjunto de los backbones modificaría la complementariedad. Tercero, las particiones repetidas de los tres datasets reutilizan las mismas imágenes y no son réplicas independientes; por ello, la inferencia se mantiene descriptiva y acompañada de tamaños de efecto y márgenes prácticos. Cuarto, CUReT opera con un desempeño cercano al techo, de modo que las diferencias entre estrategias son necesariamente pequeñas y el margen práctico, más que la magnitud observada, gobierna su interpretación. Quinto, tres datasets exploratorios fueron excluidos por problemas de procedencia o independencia, lo que reduce la variedad confirmatoria pero evita evaluaciones potencialmente contaminadas.

También existe una amenaza de generalización externa. Los tres datasets principales cubren texturas naturales, materiales y condiciones controladas, pero no representan todas las aplicaciones posibles. La composición de un subconjunto puede cambiar con la resolución, la adquisición o las clases. Por esta razón, el trabajo interpreta GFS como un procedimiento dependiente del dominio y no propone una combinación universal de extractores.

% ============================================================
% 6. CONCLUSION
% ============================================================
\section{Conclusión}\label{sec:conclusion}

Este trabajo estudió la clasificación de texturas mediante selección y concatenación de bloques de descriptores heterogéneos. El protocolo comparó, con evaluación externa separada, el mejor descriptor individual, la concatenación completa, GFS y la mejor selección homogénea. GFS mejoró el macro-F1 medio de la concatenación completa en DTD y FMD, y alcanzó equivalencia práctica en CUReT con diferencias menores que 0,001 sobre ocho biparticiones aleatorias. Utilizó entre 3,6 y 6,1 componentes en promedio y redujo la dimensionalidad entre 71,5\% y 85,2\%. Frente al mejor individual, la diferencia media fue positiva en las seis combinaciones dataset--clasificador.

La evidencia no respalda una regla según la cual toda combinación heterogénea sea superior. En una configuración de FMD y en una de CUReT, una familia homogénea obtuvo un resultado ligeramente mayor que GFS, mientras que en DTD la selección heterogénea mostró una ventaja más clara. La conclusión es que la utilidad de concatenar depende de la complementariedad efectiva y que seleccionar los bloques resulta preferible a incorporarlos indiscriminadamente cuando se busca una representación compacta.

Los controles aleatorios y la estabilidad muestran que la búsqueda suele identificar subconjuntos mejores que una elección arbitraria del mismo tamaño, pero no recupera una composición universal. En consecuencia, la contribución principal es el protocolo de selección y evaluación de bloques, no una receta fija de extractores. La evidencia disponible respalda el uso de GFS como mecanismo para obtener representaciones compactas que mejoran o mantienen el desempeño frente a concatenar indiscriminadamente todos los descriptores.

\section*{Disponibilidad de datos y código}
Los datasets utilizados son de acceso público y permanecen sujetos a las condiciones de sus respectivos distribuidores. El código experimental, los manifiestos de muestras, los hashes de procedencia, los parámetros y las tablas derivadas se organizan en el repositorio del proyecto. La URL pública y la versión archivada se incorporarán antes del envío del artículo.

\section*{Declaración ética}
El estudio utiliza exclusivamente datasets públicos de imágenes de texturas y no involucra participantes humanos, datos personales ni intervención sobre seres vivos. No se requirió aprobación de un comité de ética.

\section*{Conflicto de intereses}
Los autores declaran no tener conflictos de intereses financieros o personales que pudieran influir en este trabajo.

\section*{Financiación}
Este trabajo no recibió financiación externa específica.

\section*{Declaración sobre asistencia de herramientas de IA}
Este borrador fue preparado con asistencia de herramientas de escritura académica basadas en inteligencia artificial para apoyar la estructuración, redacción y revisión de formato. Los autores dirigieron el proceso, revisaron el contenido y asumen plena responsabilidad por su exactitud e integridad.

% --- CRediT authorship contribution statement ---
\printcredits

% --- Referencias ---
\bibliographystyle{cas-model2-names}
\bibliography{references}

\end{document}
